\section{Target Customer}
\label{sec:target-customer}

The target customer is a \textbf{Taiwan retail investor who holds or follows 5 to 20 semiconductor and technology stocks}. This user is not a professional fund manager and usually does not have access to a Bloomberg Terminal, sell-side research workflow, or a dedicated analyst team. They may use brokerage apps, finance news sites, TradingView, Goodinfo, CMoney, or spreadsheet watchlists, but their daily process is fragmented.

A concrete persona is:

\begin{quote}
    A 25--45 year-old Taiwan retail investor with a full-time job, a meaningful equity portfolio, and holdings tilted toward semiconductor, IC design, hardware, and electronics manufacturing stocks. They review holdings after work and want to know which stocks became riskier today and why.
\end{quote}

This segment is the right initial wedge for three reasons. First, the technology-stock universe is narrow enough to build a focused product instead of a generic stock platform. Second, the customer pain point is frequent and repeated: market data changes daily, while revenue and financial statements change periodically. Third, the segment already pays attention to information products such as charting tools, stock screeners, and financial data sites, which makes willingness to pay more plausible than for a purely casual investor.

The status quo is manual and scattered. A user may check brokerage app prices, search news headlines, visit separate financial statement sites, and maintain a spreadsheet. This creates three practical problems:

\begin{itemize}
    \item \textbf{Information overload}: the user sees many prices and headlines, but not a ranked list of risk changes.
    \item \textbf{Missed weak signals}: abnormal volume, drawdown, valuation pressure, and revenue deterioration may appear in different places.
    \item \textbf{No portfolio context}: most free tools focus on individual stocks rather than the user's watchlist concentration.
\end{itemize}

The proposed dashboard improves the status quo by aggregating the relevant signals into one risk-monitoring workflow. It does not replace deeper research; it prioritizes where deeper research should happen.
